\centerline{\bf \Huge Четность}

\problem{1} У каждого марсианина три руки. Могут ли семь марсиан взяться за руки?

\problem{2.1} Чётно или нечётно число 1 + 2 + 3 + ... + 2002?

\problem{2.2} Чётно или нечётно число 1 + 2 + 3 + ... + 2023?

\problem{3} Доказать: произведение
а) двух нечётных чисел нечётно;
б) чётного числа с любым целым числом чётно.

\problem{4} Разность двух целых чисел умножили на их произведение. Могло ли получиться число 2023?

\problem{5} Можно ли доску размером 5×5 заполнить доминошками размером 1×2?

\problem{6} Найдите два таких простых числа, что и их сумма, и их разность – тоже простые числа.

\problem{7} Автомат при опускании серебреника выбрасывает пять двушек, а при опускании двушки – пять серебреников. Может ли Петя, подойдя к автомату с одной двушкой, получить после нескольких опусканий одинаковое количество двушек и серебреников?

\problem{8} Два класса с одинаковым количеством учеников написали контрольную. Проверив контрольные, строгий директор Фёдор Калистратович сказал, что он поставил двоек на 13 больше, чем остальных оценок. Не ошибся ли строгий Фёдор Калистратович?

\problem{9} Даны пять чисел; сумма любых трёх из них чётна. Доказать, что все числа чётны.
  
\problem{10} На хоккейном поле лежат три шайбы А, В и С. Хоккеист бьёт по одной из них так, что она пролетает между двумя другими.
Так он делает 25 раз. Могут ли после этого шайбы оказаться на исходных местах?